\section{Global Court Query Model}
\label{sec:model}

The render pipeline supplies a camera hypothesis that cannot be learned from a
standard keypoint heatmap loss alone.  We therefore separate the model into a
global geometric encoder and deliberately lightweight dense reconstruction
heads, as shown in \cref{fig:model}.

\begin{figure}[t]
  \centering
  \resizebox{\linewidth}{!}{\begin{tikzpicture}[
  font=\sffamily\footnotesize,
  module/.style={draw, rounded corners=2pt, align=center, minimum height=9mm,
    inner xsep=4pt, inner ysep=3pt, line width=0.55pt},
  patch/.style={draw=gsblue, fill=gsblue!16, rounded corners=1pt,
    minimum width=3.2mm, minimum height=5.0mm, inner sep=0pt},
  choice/.style={draw=courtgreen!75!black, fill=white, rounded corners=1pt,
    minimum width=22mm, minimum height=4.2mm, inner sep=1pt,
    font=\sffamily\scriptsize},
  head/.style={draw=gsblue!75!black, fill=white, rounded corners=1pt,
    minimum width=20mm, minimum height=4.2mm, inner sep=1pt,
    font=\sffamily\scriptsize\bfseries},
  arrow/.style={-{Latex[length=2.1mm,width=1.35mm]}, semithick,
    draw=black!72},
  taparrow/.style={-{Latex[length=2.1mm,width=1.35mm]}, semithick,
    draw=gsblue!85!black},
  queryarrow/.style={-{Latex[length=2.1mm,width=1.35mm]}, semithick,
    draw=poseorange!90!black},
  label/.style={font=\sffamily\scriptsize, text=black!72, align=center}
]
  \node[font=\sffamily\footnotesize\bfseries, text=gsblue!85!black, anchor=west]
    at (0,5.72) {(1) Patch tokens};
  \node[font=\sffamily\footnotesize\bfseries, text=poseorange!90!black, anchor=west]
    at (7.95,5.72) {(2) Global reasoning};
  \node[font=\sffamily\footnotesize\bfseries, text=courtgreen!85!black, anchor=west]
    at (11.55,5.72) {(3) Pose + dense readouts};

  \node[module, draw=gsblue, fill=gsblue!8, minimum width=16mm,
    minimum height=14mm] (rgb) at (0.85,4.00) {};
  \fill[courtgreen!88, rounded corners=0.7pt]
    ($(rgb.center)+(-0.57,0.01)$) rectangle ($(rgb.center)+(0.57,0.47)$);
  \draw[white, line width=0.45pt]
    ($(rgb.center)+(-0.48,0.08)$) rectangle ($(rgb.center)+(0.48,0.40)$);
  \draw[white, line width=0.4pt]
    ($(rgb.center)+(0,0.08)$) -- ($(rgb.center)+(0,0.40)$)
    ($(rgb.center)+(-0.48,0.24)$) -- ($(rgb.center)+(0.48,0.24)$);
  \node[font=\sffamily\scriptsize\bfseries] at ($(rgb.center)+(0,-0.40)$)
    {RGB frame};

  \node[module, draw=gsblue, fill=gsblue!10, minimum width=18mm,
    minimum height=12mm] (dino) at (2.75,4.00)
    {{\bfseries DINOv3} ViT\\[-0.15em]\scriptsize patch tokens only};
  \node[module, draw=poseorange, fill=poseorange!9, minimum width=15mm,
    minimum height=11mm] (projection) at (4.67,4.00)
    {Patch\\projection\\[-0.2em]\scriptsize $C\!\rightarrow\!d$};

  \node[module, draw=gsblue, fill=gsblue!5, minimum width=21mm,
    minimum height=12mm] (sequence) at (6.75,4.00) {};
  \node[patch] at ($(sequence.center)+(-0.60,0.16)$) {};
  \node[patch] at ($(sequence.center)+(-0.20,0.16)$) {};
  \node[patch] at ($(sequence.center)+(0.20,0.16)$) {};
  \node[font=\sffamily\scriptsize, text=gsblue!80!black]
    at ($(sequence.center)+(0.52,0.16)$) {$\cdots$};
  \node[label] at ($(sequence.center)+(0,-0.30)$)
    {$N=H_pW_p$ projected patches};

  \node[module, draw=poseorange, fill=poseorange!10, minimum width=22mm,
    minimum height=9mm] (query) at (6.75,1.72)
    {one learned court query\\[-0.15em]\scriptsize positionless: no RoPE};

  \node[module, draw=poseorange!35, fill=poseorange!3, minimum width=24mm,
    minimum height=14mm] at (9.58,4.18) {};
  \node[module, draw=poseorange!55, fill=poseorange!5, minimum width=24mm,
    minimum height=14mm] at (9.48,4.09) {};
  \node[module, draw=poseorange, fill=poseorange!11, minimum width=24mm,
    minimum height=14mm] (task) at (9.38,4.00)
    {Global task Transformer\\[-0.10em]
     \scriptsize shared self-attention + MLP\\[-0.10em]
     \scriptsize 2D RoPE on patch $Q,K$ only};

  \draw[arrow] (rgb.east) -- (dino.west);
  \draw[arrow] (dino.east) -- (projection.west);
  \draw[arrow] (projection.east) -- (sequence.west);
  \draw[arrow] (sequence.east) -- (task.west);
  \draw[queryarrow] (query.east) -| ($(task.south west)+(0.28,0)$);

  \node[module, draw=poseorange, fill=poseorange!10, text width=18mm,
    minimum height=9mm] (posehead) at (12.10,4.52)
    {final query\\[-0.15em]\scriptsize MLP pose head};
  \node[module, draw=poseorange, fill=poseorange!6, text width=32mm,
    minimum height=12mm] (poseout) at (15.15,4.52)
    {\bfseries raw pose10D\\[-0.10em]
     \scriptsize camera centre $\hat{\mathbf C}$ (3)\\[-0.10em]
     \scriptsize row--6D $\hat{\mathbf a}$ (6) $\;\vert\;$
     $\widehat{\log f}$ (1)};

  \node[choice] (linear) at (12.18,2.18) {Linear: $1{\times}1$ + resize};
  \node[choice] (progressive) at (12.18,1.70) {Progressive: upsample + DW conv};
  \node[choice] (dpt) at (12.18,1.22) {DPT: reassemble + fuse taps};
  \begin{scope}[on background layer]
    \node[draw=courtgreen, fill=courtgreen!8, rounded corners=2pt,
      fit=(linear)(progressive)(dpt), inner xsep=2.5mm, inner ysep=2.2mm]
      (decoders) {};
  \end{scope}
  \node[font=\sffamily\scriptsize\bfseries, text=courtgreen!85!black]
    at ($(decoders.north)+(0,0.23)$) {lightweight decoder (choose one)};

  \node[head] (kphead) at (15.28,2.18) {KP14};
  \node[head] (linehead) at (15.28,1.70) {LINE};
  \node[head] (seghead) at (15.28,1.22) {SEG};
  \begin{scope}[on background layer]
    \node[draw=gsblue, fill=gsblue!7, rounded corners=2pt,
      fit=(kphead)(linehead)(seghead), inner xsep=2.5mm, inner ysep=2.2mm]
      (heads) {};
  \end{scope}
  \node[font=\sffamily\scriptsize\bfseries, text=gsblue!85!black]
    at ($(heads.north)+(0,0.23)$) {selected $1{\times}1$ heads};

  \draw[queryarrow] (task.east) |- (posehead.west);
  \draw[queryarrow] (posehead.east) -- (poseout.west);
  \draw[taparrow] ($(task.south east)+(-0.23,0)$) |- (decoders.west);
  \node[label, text=gsblue!85!black, anchor=east] at (10.32,2.82)
    {task-layer patch taps};
  \draw[taparrow] (decoders.east) -- (heads.west);

  \node[label, text width=5.8cm] at (13.75,0.35)
    {Dense heads restore image space; global mixing stays in the task Transformer.};
\end{tikzpicture}
}
  \caption{\textbf{Global query, mechanical dense decoding.}  \DINO{} patch
  tokens are projected and receive 2D rotary encoding.  One learned court query
  receives no positional rotation but globally attends with every patch in the
  task Transformer.  Its final state predicts pose10D; selected patch taps feed
  interchangeable lightweight decoders and direct \KP, line, and segmentation
  heads.}
  \label{fig:model}
\end{figure}

\subsection{Strict patch-token boundary}

Given an augmented image \(I\in\R^{H\times W\times 3}\), we replicate-pad its
right and bottom boundaries to the \DINO{} patch multiple.  The backbone
returns a sequence containing patch and special tokens; the adapter accepts
only the \(H_pW_p\) patch tokens and verifies that token count, patch grid,
dtype, device, and finiteness agree.  CLS and register tokens are excluded so
that the task encoder has one unambiguous global token: its own learned court
query.

The patch features \(\mathbf{x}_{ij}\) are projected to task width \(d\).  For
each attention layer, 2D rotary position encoding rotates the patch query/key
pairs as a function of row and column \citep{su2021roformer}.  The prepended
learned token \(\mathbf{q}_0\) is not rotated:
\begin{equation}
  \mathbf{Z}^{(0)}=
  [\mathbf{q}_0,\,\mathbf{P}\mathbf{x}_{00},\ldots,
   \mathbf{P}\mathbf{x}_{H_p-1,W_p-1}],
  \qquad
  \operatorname{RoPE}_{2D}(\mathbf{q}_0)=\mathbf{q}_0.
  \label{eq:query_input}
\end{equation}
All query and patch tokens participate in the same self-attention layers.
Thus the final query state \(\mathbf{q}_L\) can aggregate any visible part of
the court, while patch states retain their spatial identity.

\subsection{One camera hypothesis}
\label{sec:pose10d}

The target view is first canonicalized by a proper camera-end rotation
\begin{equation}
  \T{\Canon}{\Cam}=\mathbf{S}\T{\Court_i}{\Cam},
  \qquad \mathbf{S}\in\{\mathbf{I},\mathbf{R}_z(\pi)\}.
  \label{eq:canonical_pose}
\end{equation}
The physical \KP ordering and canonicalization choice are stored authority;
the model adapter rejects a disagreement between the dense target and pose
target rather than inferring a permutation.

An MLP maps \(\mathbf{q}_L\) to
\begin{equation}
  \widehat{\mathbf{p}}=
  [\widehat C_x,\widehat C_y,\widehat C_z,
   \widehat a_{11},\widehat a_{12},\widehat a_{13},
   \widehat a_{21},\widehat a_{22},\widehat a_{23},
   \widehat\ell_f]\in\R^{10}.
  \label{eq:pose10d}
\end{equation}
The first three values are the camera centre \(\widehat{\mathbf{C}}\) in
canonical court metres.  The next two raw row vectors are orthogonalized:
\begin{align}
  \mathbf{r}_1&=\frac{\mathbf{a}_1}{\|\mathbf{a}_1\|_2}, &
  \bar{\mathbf{r}}_2&=\mathbf{a}_2-
      (\mathbf{r}_1^{\top}\mathbf{a}_2)\mathbf{r}_1, \\
  \mathbf{r}_2&=\frac{\bar{\mathbf{r}}_2}{\|\bar{\mathbf{r}}_2\|_2}, &
  \mathbf{r}_3&=\mathbf{r}_1\times\mathbf{r}_2, \\
  \widehat{\mathbf{R}}_{\Canon\leftarrow\Cam}
    &= [\mathbf{r}_1^{\top};\mathbf{r}_2^{\top};\mathbf{r}_3^{\top}], &
  \widehat f&=\exp(\widehat\ell_f).
  \label{eq:rotation_decode}
\end{align}
This continuous representation avoids quaternion normalization/sign ambiguity
and is rejected if either basis vector degenerates
\citep{zhou2019continuity}.  Focal length is finite and strictly positive; the
principal point is not predicted but comes from the augmentation-aware target
intrinsics.

\subsection{Lightweight dense heads}

Selected patch states from the task Transformer feed one of three controlled
decoder families:
\begin{itemize}
  \item \emph{linear}: a per-location projection followed by resize;
  \item \emph{progressive}: staged upsampling with depthwise-separable
    refinement; and
  \item \emph{DPT}: multi-tap reassembly and feature fusion.
\end{itemize}
The decoder produces an image-resolution feature map.  Independent \(1\times1\)
heads emit fourteen singleton keypoint logits, one line-mask logit map, and
court segmentation logits.  Decoder width, stages, fusion mode, and taps are
typed ablation variables, but the output contract never changes.  This design
tests whether global geometry can reside in the shared token encoder instead of
being reconstructed by an opaque high-capacity task head.

The direct objective remains indispensable:
\begin{align}
  \mathcal{L}_{\mathrm{dense}}
   &=w_{\mathrm{kp}}\mathcal{L}_{\mathrm{focalBCE}}
    +w_{\mathrm{line}}\mathcal{L}_{\mathrm{BCE+Dice}}
    +w_{\mathrm{seg}}\mathcal{L}_{\mathrm{CE+Dice}}, \\
  \mathcal{L}_{\mathrm{pose}}
   &=w_t\|\widehat{\mathbf{C}}-\mathbf{C}\|_2^2
    +w_R\,d_{\SOthree}(\widehat{\mathbf{R}},\mathbf{R})
    +w_f(\widehat\ell_f-\ell_f)^2,
  \label{eq:direct_losses}
\end{align}
where \(d_{\SOthree}\) is geodesic rotation distance.  The auxiliary objective
next connects these otherwise independent branches; it never replaces their
ground-truth losses.

